\documentclass[11pt]{article}
\usepackage{preprint}

\title{Political Institutions and Economic Outcomes: A Quantitative Analysis of US States}
\author{Jiadi He\email{jh2288@georgetown.edu}}
\affil{Georgetown University}
\shorttitle{Political Institutions and Economic Outcomes}
\keywords{political data; institutions; economic development}

\begin{document}
\maketitle
\keywordsline

\section{Introduction}

The primary role of political research is to describe various political concepts and analyze the relationships between them. The dataset used in this research is the States dataset for R Companion to Essentials of Political Analysis. The dataset contains information on demographics, economics, and political variables in the fifty US states. The politics of states significantly influence the state's economic development. Several studies have shown that politics influence the economics of states and regions in a country. The primary objective of this research is to investigate the relationship between political affiliations of state leadership or affiliation and economic developments in the region. The research question is; what is the influence of a state's political affiliation on the economic performance of the state? This research question is important in addressing the economic development imbalance witnessed in areas where regional leaders are not affiliated with the ruling party.

\section{Theories}

A study by Asher \& Novosad (2017) indicated that political favoritism affects the allocation of government resources and, thus, economic development. It is well-established in the scholarly literature that both wealthy and developing nations' political systems affect the timing and distribution of government interventions. While other factors such as availability of natural resources, infrastructure, and tax policies, influence economic growth and development, political leadership has an equal influence on the economic development of a region.

The strength of a political party is also a significant determinant in the economic progression of a country. Economic performance may be affected by the power of political parties, according to studies of both democracies and autocracies. Studies have shown that the influence of political parties on economic growth may hinge on factors inside each party and those between them. A political party's capacity to regulate the economy is influenced by its structure (Bizzarro et al., 2018).

\section{Hypotheses}

The null hypothesis drawn from the research question states that the party affiliation of state leadership, including the governor, has no significant relationship on the state's economic performance. The alternative would be that the party affiliation of state leadership, including the governor, significantly links to the state's economic performance.

\section{Data, Methods, and Variables}

As stated earlier, the data for this research will be drawn from the States dataset for R Companion to Essentials of Political Analysis. The dataset contains information on demographics, economics, and political variables in the fifty US states. Therefore, this data can infer the relationship between the political affiliations of state leadership and the state's economic performance. The dependent variable, also called the outcome variable, in this research is the state's economic performance. The independent variable is the state leadership's party or political affiliation, including that of the governor and the legislators. The independent variable predicts or explains the outcome variable (Jenkins-Smith et al., 2017). Therefore, in this case, political affiliation is used to infer the outcome variable: a state's economic performance. A state's economic performance is measured by the real gross domestic product (GDP). GDP is a valid and reliable method of measuring a country's economic performance. The term “validity” describes how well the measurement instrument represents the concept. In contrast, reliability describes how well the measurement holds up over time.

The dependent variable, economic performance, can be described as a numerical variable. On the other hand, the independent variable (political affiliation) can be described as a categorical variable. A politician in the United States is either a democrat or a republican. The study will use regression analysis to test the hypothesis. As a method for modeling the connections between the explanatory factors and the modeled dependent variable, regression analysis will be used in this study. The regression coefficient is a statistic calculated from a regression analysis that provides a rough approximation of the relationship of an independent variable on a dependent one (Pollock et al., 2022). That way, the study will establish the relationship of political affiliation on the states' economic performance.

\section{Descriptive Statistics}

The analysis of descriptive statistics is the most basic form of analysis that can be performed in research. Measures of central tendency, such as the mean, median, and mode, will be evaluated in the study. Political affiliation, a categorical nominal variable, categorizes states' political leadership into either democrats or republicans. In R, nominal variables like this are described as unordered factors (Pollock et al., 2022). For that reason, one can determine the mode, which describes the most common value of the variable. That way, the study can establish the most common political affiliation in the 50 states. These values can be tabulated in a frequency distribution table. At the same time, the strength of the political parties in the state, which is an ordinal-level variable, can be gauged. R refers to ordinal variables as ordered factors. All measures of central tendency, including the mean, median, and mode, can be evaluated for the dependent variable. This research will also use bar charts drawn from frequency distribution tables. The distribution of the dependent and the dependent variables across the 50 states will be drawn from the dataset used in the study.

\section{Hypothesis Testing and Interpretations}

The preliminary results suggest that political affiliation has a significant association with the economic performance of states. Therefore, the research rejects the null hypothesis and accepts the alternative one. These findings correlate with findings from similar other studies. Studies contend that Democratic governments generate faster economic growth than their republican counterparts. At the state level, ideology-induced policies are more common, influencing the state's economic performance. Compared to their conservative counterparts, studies suggest that Democratic governors are more likely to enact comprehensive and liberal programs. Contrary to what may be believed at the national level, local government ideology has little to no influence on economic policy (Potrafke, 2018). These findings indicate that the political ideologies of the parties significantly influence the states' economic performance.

\section{Limitations}

The research does have certain limitations. While the study has evaluated the relationship of political affiliation on states' economic performance, it is crucial to consider that there are more variables other than political affiliation that affect economic performance. Cross-sectional data, which is what this study has, makes it difficult to draw conclusions about the relationship between political affiliation and economic performance. The study is also limited in that other characteristics, such as political philosophy or personal policy preferences, are not considered.

Lastly, other indicators of economic growth, such as unemployment or poverty rates, are not considered in the research, which focuses solely on GDP per capita. Therefore, there is a need to evaluate other socio-economic and demographic factors that influence the performance of both Democratic and Republican states. It is also important to evaluate how the implementation of policies in democratic and republican states influences economic performance.

\section{Conclusion}

This paper has evaluated the relationship between state leadership's political affiliation, such as governors, and the state's economic performance. The study used the States dataset for R Companion to Essentials of Political Analysis, which contains information on demographics, economics, and political variables in the fifty US states. Several studies have concluded that political parties matter in the United States. Numerous researchers have examined the relationship of politicians' and policymakers' ideological leanings and party affiliations on economic policy-making. Partisan theories propose that Democratic governments are more likely to enact expansive measures than their Republican counterparts (Potrafke, 2018). Democratic states, given the moniker blue states, control the highest percentage of the country's GDP (Green, 2020). The findings of this research have implications for understanding the influence of ideologies fronted by political parties in the country and how they influence economic performance at the state and federal levels.

\begin{references}
\item Asher, S., \& Novosad, P. (2017). Politics and local economic growth: Evidence from India. \textit{American Economic Journal: Applied Economics}, \textit{9}(1), 229-273.
\item Bizzarro, F., Gerring, J., Knutsen, C. H., Hicken, A., Bernhard, M., Skaaning, S. E., ... \& Lindberg, S. I. (2018). Party strength and economic growth. \textit{World Politics}, \textit{70}(2), 275-320.
\item Green, A. (2020). \textit{The Growing Divide: Red States vs. Blue States - Georgetown Public Policy Review}. Georgetown Public Policy Review. \url{https://gppreview.com/2020/02/21/growing-divide-red-states-vs-blue-states/}
\item Jenkins-Smith, H., Ripberger, J. T., Copeland, G., Nowlin, M. C., Hughes, T., Fister, A. L., \& Wehde, W. (2017). \textit{Quantitative Research Methods for Political Science, Public Policy and Public Administration (With Applications in R)}. University of Oklahoma Libraries.
\item Pollock III, P. H., \& Edwards, B. C. (2022). \textit{An R companion to political analysis}. CQ Press.
\item Potrafke, N. (2018). Government ideology and economic policy-making in the United States—a survey. \textit{Public Choice}, \textit{174}(1-2), 145-207.
\end{references}

\end{document}
